% ===================================================================
%  तक्ता क्रमांक १६ — मोठे दुकान. --- बाजार. खेळणी.
%
%  Traduction, faite en 2026, en marathi d'aujourd'hui, sur l'ido de
%  texto/io. Regles de la colonne : voir 10-takta-01.tex. Dernier
%  tableau du livret. Une seule suite d'alineas numerotes : les cles
%  ne portent pas d'indice de scene. Le fac-simile compose « d) »
%  sans sa parenthese ouvrante ; la traduction ecrit la reference
%  bien formee, comme au tableau 8.
%
%  LE CIEL EST LE MEILLEUR MORCEAU DE CE DERNIER TABLEAU, ET C'EST
%  UNE SURPRISE. Le grand-pere s'arrete devant un tableau
%  d'astronomie que l'enfant ne comprend pas ; le marathi, lui, le
%  comprend entierement, parce que chacun de ces mots est en sanskrit
%  et vieux de deux mille ans : धूमकेतू la comete — celle qui porte
%  une banniere de fumee —, ग्रहण l'eclipse, ग्रह les planetes,
%  गोलार्ध les hemispheres, सूर्यमाला le systeme solaire, et les
%  quatre points cardinaux उत्तर, दक्षिण, पूर्व, पश्चिम. Le seul
%  passage du livret que l'enfant declare ne pas comprendre est celui
%  ou la colonne n'a rien eu a emprunter.
%
%  ET LA GUERRE DE JOUET A DES GRADES MARATHES : सुभेदार le
%  sergent-major, हवालदार le sergent d'artillerie, जायबंदी
%  l'invalide, मोहीम la campagne, रणभूमी le champ de bataille,
%  पराभव la defaite, तह le traite. Les Marathes ont fait la guerre
%  assez longtemps pour en nommer chaque piece.
%
%  LA MENAGERIE EN CARTON EST ENTIEREMENT MARATHE — पोपट le
%  perroquet, गेंडा le rhinoceros, शहामृग l'autruche, माकड le singe,
%  कोल्हा le renard, सुसर le crocodile, उंट le chameau, बिबट्या le
%  leopard, तरस la hyene, लांडगा le loup, पाणघोडा l'hippopotame,
%  वटवाघूळ la chauve-souris, अजगर le boa, कासव la tortue, देवमासा
%  la baleine — le « poisson-dieu » —, मुशी le requin. Ne restent
%  etrangers que le renne, la belette et le dromadaire.
%
%  ET LA BARAKHADI EST LE DERNIER MOT QUE LE MARATHI AVAIT DEJA.
%  L'abecedaire (80) s'appelle बाराखडी, les douze mesures : les douze
%  voyelles qu'on decline apres chaque consonne pour apprendre a
%  lire le devanagari. Le livret dit « alfabet-libro », qui compte
%  des lettres ; le marathi compte des syllabes, parce que son
%  ecriture en est faite. C'est la derniere chose que ce tableau
%  apprend, et elle est dans le nom meme du livre ou l'on apprend.
% ===================================================================

%%K t16-apar-1 apar
\VUsaut{4.0mm}
\VUtitre{15.0pt}{60}{तक्ता क्रमांक १६}
\VUsaut{2.0mm}
\VUpk{13.2pt}{40}{मोठे दुकान. --- बाजार.}
\VUpk{13.2pt}{40}{खेळणी.}
\VUsaut{3.4mm}

%%K t16-01-1 p
१. --- नवीन वर्ष जवळ येते आहे. मोठ्या दुकानांनी \VUgras{खेळण्यांच्या} \textsuperscript{(1)}
आणि नववर्षाच्या भेटींच्या आपल्या मांडण्या तयार केल्या आहेत.

%%K t16-01-2 p
मी माझ्या आजोबांना विनंती केली की त्यांनी मला ते सुंदर प्रदर्शन पाहायला बाजारात न्यावे, आणि
त्यांनी होकार दिला.

%%K t16-02-1 p
२. --- आधी आम्ही सुंदर शस्त्रांच्या फलकांसमोर थांबलो, ज्यांच्यावर \VUgras{बंदूक},
\VUgras{टोपी}, \VUgras{समशेर}, \VUgras{तलवारीचा पट्टा} आणि \VUgras{खांदेपट्ट्यांची} जोडी होती,
किंवा \VUgras{शिरस्त्राण}, \VUgras{चिलखत} \textsuperscript{(3)} आणि \VUgras{पिस्तूल}
\textsuperscript{(4)}. हे सगळे मला \VUgras{प्रकाशझोत यंत्रांपेक्षा} \textsuperscript{(2)}
कितीतरी जास्त रोचक वाटले.

%%K t16-tit-1 sub
\VUsaut{3.0mm}
\VUpk{10.2pt}{40}{सुंदर देखावे.}
\VUsaut{2.6mm}

%%K t16-03-1 p
३. --- त्याच विभागात आपल्या \VUgras{खोपटीत} एक \VUgras{पहारेकरी} \textsuperscript{(5)} होता;
तो पुठ्ठ्याच्या एका अख्ख्या प्राणिसंग्रहासमोर पहारा देत असल्यासारखा दिसत होता. तिथे केवढे
प्राणी होते : आपल्या दांडीवर बसलेला \VUgras{पोपट} \textsuperscript{(6)}, नाकावर शिंग असलेला
\VUgras{गेंडा} \textsuperscript{(7)}, \VUgras{रेनडियर} \textsuperscript{(8)}, \VUgras{शहामृग}
\textsuperscript{(9)}, अक्रोड फोडणारे \VUgras{माकड} \textsuperscript{(10)}, कोंबडी पळवणारा
\VUgras{कोल्हा} \textsuperscript{(11)}, अवाढव्य जबडा उघडणारी \VUgras{सुसर}
\textsuperscript{(12)}, \VUgras{उंट} \textsuperscript{(13)}, डौलदार \VUgras{ग्रेहाउंड}
\textsuperscript{(14)}, \VUgras{पँथर} \textsuperscript{(15)}, आणि मग मला माहीत नसलेले दोन
प्राणी, जे म्हणे \VUgras{वीझल} \textsuperscript{(16)} आणि \VUgras{तरस} \textsuperscript{(17)}
आहेत. तिथे \VUgras{बिबट्याही} \textsuperscript{(18)} होता आणि \VUgras{लांडगाही}
\textsuperscript{(21)}; अगदी जवळ रबराच्या नाजूक \VUgras{घुशी} \textsuperscript{(19)} आणि
इवलीशी \VUgras{उंदरे} \textsuperscript{(20)} होती. शेवटी \VUgras{पाणघोड्याने}
\textsuperscript{(22)} आणि वशिंड असलेल्या \VUgras{ड्रोमेडरीने} \textsuperscript{(23)} तो
संग्रह पुरा केला. शेजारीच \VUgras{शिशाच्या शिपायांनी} \textsuperscript{(29)} भरलेली भव्य छावणी
नसती तर मी तिथे आणखी कितीतरी तास राहिलो असतो — तिनेच माझे लक्ष वेधून घेतले.

%%K t16-tit-2 sub
\VUsaut{4.0mm}
\VUpk{10.2pt}{40}{शिपाई आणि युद्ध.}
\VUsaut{2.8mm}

%%K t16-04-1 p
४. --- माझे आजोबा \VUgras{जायबंदी} \textsuperscript{(24)} आहेत आणि एका \VUgras{जखमेमुळे}
त्यांना सैन्य सोडावे लागले — त्याच जखमेमुळे त्यांना \VUgras{लष्करी पदक} मिळाले —, आणि ज्या
\VUgras{मोहिमांत} ते सामील झाले त्यांच्या गोष्टी मला सांगायला त्यांना फार आवडते. त्या इवल्याशा
सैन्याच्या निरनिराळ्या हालचाली समजावून सांगताना ते सुखावले. बाजार पाहायला आलेल्या खऱ्या
शिपायांचा एक घोळका — एक \VUgras{अभियंता-शिपाई} \textsuperscript{(25)}, \VUgras{बेरे} घातलेला
एक \VUgras{आल्प्सचा शिकारी} \textsuperscript{(26)}, पायदळाचा एक \VUgras{सुभेदार}
\textsuperscript{(27)} आणि बाहीवर सोनेरी \VUgras{पट्टी} असलेला तोफखान्याचा एक \VUgras{हवालदार}
\textsuperscript{(28)} — ते स्पष्टीकरण ऐकायला आमच्याजवळ थांबला.

%%K t16-05-1 p
५. --- एवढे झळाळते श्रोते मिळाल्यामुळे फार अभिमानाने माझ्या आजोबांनी तिथल्या तिथे एक पुरी
लढाईची योजना रचली.

%%K t16-05-2 p
आपले \VUgras{तट} आणि \VUgras{खंदक} असलेल्या तटबंदीच्या शहराला \VUgras{शत्रुसैन्याने} वेढा
घातला होता, आणि दीर्घ \VUgras{तोफांच्या माऱ्यानंतर} ते \VUgras{हल्ल्याला} तयार होते. पण भुकेने
पिचलेल्या \VUgras{वेढलेल्यांनी} \VUgras{वेढा घालणाऱ्यांच्या} \VUgras{छावणीवर} \VUgras{चढाई}
करून पाहिली. \VUgras{तंबूंच्या तळाभोवती} चिवट \VUgras{लढाई} देऊन \VUgras{हल्ला} परतवल्यावर
\VUgras{जिंकलेल्या} वेढेकऱ्यांनी \VUgras{हरलेल्या} हल्लेखोरांचा पाठलाग करायला आपल्या
\VUgras{घोडदळाच्या तुकड्या} सोडल्या. त्यांनी \VUgras{माघार} घेतली, \VUgras{रणभूमी}
\VUgras{मेलेल्यांनी} आणि \VUgras{जखमींनी} भरून टाकत. \VUgras{स्ट्रेचरवाल्यांनी} या शेवटच्यांना
\VUgras{रुग्णवाहिकेत} नेले, \VUgras{शल्यचिकित्सकाकडून} उपचार करून घ्यायला.

%%K t16-05-3 p
\VUgras{चौरस} रचना केलेल्या काही \VUgras{पलटणींनी} वीरतेने प्रतिकार केला तरी \VUgras{पराभव} पुरा झाला, उरलेले सैन्य \VUgras{पळाले}, पुष्कळ \VUgras{कैदी} मागे सोडून. शत्रू शहर ताब्यात घेऊन आपला \VUgras{विजय} पुरा करायला निघाला. कदाचित हा मोहिमेचा शेवट असेल, आणि \VUgras{युद्धविरामानंतर} \VUgras{शांततेचा तह} करता येईल.

%%K t16-tit-3 sub
\VUsaut{4.4mm}
\VUpk{10.2pt}{40}{नववर्षाच्या भेटी.}
\VUsaut{2.8mm}

%%K t16-06-1 p
६. --- त्या अनुभवी वीराचे चित्रमय वर्णन ऐकून हसणाऱ्या तरुण शिपायांनी त्यांच्याकडे आणखी काही
खुलासे मागितले. मी मात्र सुंदर \VUgras{क्रॉसबो} \textsuperscript{(32)} आणि \VUgras{धनुष्य}
\textsuperscript{(33)} आपल्या \VUgras{बाणासह} \textsuperscript{(31)} पाहायला दुकानभर फिरलो.
तिथे मला प्राण्यांचा आणखी एक संग्रह सापडला : दोन \VUgras{गाणारे पक्षी} \textsuperscript{(34)}
— यांत्रिक \VUgras{नाइटिंगेल} आणि पूर्ण गळ्याने गाणारा \VUgras{वटवट्या} —, आणि विचित्र
प्राण्यांची एक मालिका : \VUgras{वटवाघूळ} \textsuperscript{(35)}, \VUgras{अजगर}
\textsuperscript{(36)}, \VUgras{कासव} \textsuperscript{(37)}, \VUgras{सील}
\textsuperscript{(38)}, \VUgras{देवमासा} \textsuperscript{(39)} आणि खोटी \VUgras{मुशी}
\textsuperscript{(40)}.

%%K t16-07-1 p
७. --- ते न्याहाळायला मी फार वेळ थांबलो नाही, कारण मला माझ्या स्वप्नातली गोष्ट दिसली : भव्य
सजावट आणि मोहक लाल पडदा असलेले एक अतिशय सुंदर \VUgras{कठपुतळ्यांचे नाट्यगृह}
\textsuperscript{(41)}. \VUgras{रंगमंचावर} दोन नट एका फार रोचक \VUgras{नाटकात} \VUgras{भूमिका}
करत असल्यासारखे दिसत होते, कारण कक्षांत बसवलेले किंवा \VUgras{आरामखुर्च्यांवर} आणि
\VUgras{वाद्यवृंदाच्या प्रमुखामागच्या} \VUgras{खालच्या भागात} बसलेले पुठ्ठ्याचे इवलेसे
\VUgras{प्रेक्षक} जोरात टाळ्या वाजवत होते.

%%K t16-07-2 p
दुर्दैवाने ते नाट्यगृह फार महाग होते.

%%K t16-08-1 p
८. --- शिवाय, खेळणी निवडणे अवघड होते, कारण दर्शनी काचांमध्ये सगळ्या प्रकारची खेळणी साचली होती
: \VUgras{आर्लेकँ} \textsuperscript{(42)}, \VUgras{पुल्चिनेल्ला} \textsuperscript{(43)},
\VUgras{पिएरो} \textsuperscript{(44)}, पंख फडकावणारी \VUgras{घुबडे} \textsuperscript{(45)},
शीळ घालणारी \VUgras{काळी कस्तूर} \textsuperscript{(46)}, भुंकणारे \VUgras{पूडल},
\VUgras{ग्रामोफोन} \textsuperscript{(47)} आणि \VUgras{संयमाची कोडी}. यांतल्या एका खेळण्याने
माझी फार करमणूक केली : ते एक \VUgras{नाविक लढाई} \textsuperscript{(48)} दाखवत होते, जिच्यात
\VUgras{नारळाच्या झाडांनी} भरलेल्या एका किनाऱ्याजवळ \VUgras{चाच्यांनी} एका \VUgras{नौकेवर}
हल्ला केला होता.

%%K t16-noto-1 noto
\VUnotes{143.2mm}{%
\textit{(*) तक्ता क्रमांक ६ पाहा.}%
}

%%K t16-09-1 p
९. --- ती सगळी आश्चर्ये मला फार सहजपणे न्याहाळता आली, कारण दुकानात थोडीच माणसे होती — जेमतेम
काही रिकामटेकडे, एक \VUgras{ड्रॅगून} \textsuperscript{(49)} आणि एक \VUgras{वसुली करणारा}
\textsuperscript{(50)}, जो \VUgras{मुख्य गल्ल्यावर} \textsuperscript{(51)} \VUgras{हुंडी} सादर
करत होता.

%%K t16-10-1 p
१०. --- \VUgras{गल्लेवालीच्या} मागे फळ्यांवर ठेवलेल्या पुस्तकांचा उपयोग काय असे मी माझ्या
आजोबांना विचारले; त्यांनी मला सांगितले की त्या \VUgras{हिशेबाच्या वह्या} आहेत.

%%K t16-tit-4 sub
\VUsaut{3.6mm}
\VUpk{10.2pt}{40}{दुकानाभोवती टकमक.}
\VUsaut{2.8mm}

%%K t16-11-1 p
११. --- खेळण्यांचा विभाग सोडून आम्ही \VUgras{खोगिरे} \textsuperscript{(53)}, \VUgras{रिकिबी}
\textsuperscript{(54)}, \VUgras{लगाम} \textsuperscript{(55)}, \VUgras{तोंडातल्या कड्या}
\textsuperscript{(56)} आणि \VUgras{चाबूक} पाहायला खोगीरविभागाकडे गेलो. \VUgras{विभागप्रमुख}
\textsuperscript{(57)} माझ्या आजोबांना काही माहिती देत असताना मी \VUgras{पोर्सिलेन} आणि
\VUgras{स्फटिकाची} \textsuperscript{(58)} मांडणी पाहायला गेलो, जिथे सुंदर
\VUgras{कोशिंबिरीची भांडी} आणि नाजूक \VUgras{चहाच्या किटल्या} \textsuperscript{(59)} आहेत.

%%K t16-12-1 p
१२. --- पहिल्या मजल्यावर जायला आम्ही \VUgras{कॉर्सेटच्या} \textsuperscript{(60)} दर्शनी
काचेमागून गेलो, मोजांच्या दुकानाशेजारून, जिथे फार सुंदर \VUgras{आतले कपडे}
\textsuperscript{(63)} मांडले होते. जवळच एक \VUgras{विक्रेतीण} \textsuperscript{(61)} एका
\VUgras{गिऱ्हाईकबाईकडून} \textsuperscript{(62)} सुंदर रेशमी \VUgras{कापडे}
\textsuperscript{(64)} निवडून घेत होती.

%%K t16-12-2 p
जिन्याने वर चढताना माझ्या लक्षात आले की दुकान जपानी \VUgras{कंदिलांनी} \textsuperscript{(65)},
\VUgras{सूर्यछत्र्यांनी} \textsuperscript{(66)} आणि अवाढव्य \VUgras{ताडांनी}
\textsuperscript{(70)} सजवले आहे.

%%K t16-13-1 p
१३. --- पहिल्या मजल्यावर पाहण्यासारखे फार काही नव्हते; फक्त फर्निचर (*), \VUgras{तिजोऱ्या}
\textsuperscript{(67)}, \VUgras{खणांची कपाटे} \textsuperscript{(68)} आणि \VUgras{बाळगाड्या}
\textsuperscript{(69)}. पण दुकानाचे एकंदर दृश्य फार \VUgras{मजेदार} होते.

%%K t16-14-1 p
१४. --- आमच्या पायांशी मला वाद्ये दिसली : \VUgras{हार्प} \textsuperscript{(71)},
\VUgras{गिटार} \textsuperscript{(72)}, \VUgras{मॅन्डोलिन} \textsuperscript{(73)},
\VUgras{व्हॉल्व्हची कर्णे} \textsuperscript{(74)}, \VUgras{क्लॅरिनेट} \textsuperscript{(75)},
आपल्या \VUgras{गजासह} \textsuperscript{(76)} व्हायोलिन आणि \VUgras{नगारासुद्धा}
\textsuperscript{(77)}. जरा पुढे, कागदाच्या विभागाजवळ, एक \VUgras{विद्यार्थी}
\textsuperscript{(78)} \VUgras{विक्रेत्याकडून} \textsuperscript{(79)} वह्यांच्या पिशवीचा भाव
करत होता. त्यांच्याजवळ मला एक सुंदर \VUgras{बाराखडीचे पुस्तक} \textsuperscript{(80)} दिसले.

%%K t16-14-2 p
दुकानाच्या मधोमध एक उंच \VUgras{नाताळचे झाड} \textsuperscript{(81)} उभे केले होते,
मेणबत्त्यांनी, लहानसहान वस्तूंनी आणि सगळ्या प्रकारच्या खेळण्यांनी पुरे भरलेले :
\VUgras{लाकडी घोडे} \textsuperscript{(82)}, \VUgras{बाहुल्या} \textsuperscript{(83)}, वगैरे.

%%K t16-14-3 p
आपली \VUgras{दंतमंजने} आणि निरनिराळी \VUgras{अत्तरे} असलेल्या
\VUgras{सुगंधद्रव्यांच्या दुकानातून} \textsuperscript{(84)} फार चांगला वास पसरत होता, जो
आमच्यापर्यंत येत होता.

%%K t16-tit-5 sub
\VUsaut{3.4mm}
\VUpk{10.2pt}{40}{आपण काहीतरी विकत घेतो.}
\VUsaut{2.8mm}

%%K t16-15-1 p
१५. --- माझ्या आजोबांना वेळ फार लांबल्यासारखा वाटू लागल्यामुळे आम्ही मोठ्या जिन्याने खाली
उतरलो. \VUgras{खगोलशास्त्राच्या तक्त्यासमोर} \textsuperscript{(85)} ते थोडे थांबले, जो — ते
मला म्हणाले — \VUgras{धूमकेतू} \textsuperscript{(a)}, \VUgras{चंद्र आणि सूर्यग्रहणे}
\textsuperscript{(b)}, \VUgras{चार मुख्य दिशा} \textsuperscript{(c)}
\VUgras{(उत्तर, दक्षिण, पूर्व आणि पश्चिम)} दाखवणारे वाऱ्याचे चक्र, दोन \VUgras{गोलार्धांचा}
\textsuperscript{(d)} नकाशा, \VUgras{ग्रह} \textsuperscript{(e)} आणि \VUgras{सूर्यमाला}
\textsuperscript{(f)} दाखवतो. यातले मला काहीच कळले नाही, आणि तिथून गेलेल्या एका
\VUgras{पोहोचवणाऱ्या नोकराच्या} \textsuperscript{(86)} पट्ट्या लावलेल्या गणवेशात, आणि
माझ्याजवळ \VUgras{पैशाच्या पिशव्यांशेजारी} \textsuperscript{(88)} व
\VUgras{कागदपत्रांच्या पिशव्यांशेजारी} \textsuperscript{(89)} असलेल्या सुंदर \VUgras{पंख्यांत}
\textsuperscript{(87)} मला कितीतरी जास्त रस होता.

%%K t16-16-1 p
१६. --- मांडणीवरची \VUgras{पिसेही} \textsuperscript{(90)} आणि \VUgras{पट्ट्यांची बकले}
\textsuperscript{(91)} मला आवडली, आणि टोपल्यांत डौलदारपणे मांडलेली ती सुंदर
\VUgras{कृत्रिम फुलेही} \textsuperscript{(94)}. \VUgras{वायलेट}, \VUgras{जिलीफ्लॉवर},
\VUgras{हायसिंथ} \textsuperscript{(95)}, \VUgras{जिरॅनियम} \textsuperscript{(96)},
\VUgras{लिली} \textsuperscript{(97)} आणि \VUgras{नास्टर्शियम} \textsuperscript{(98)} फार
चांगले उतरवले होते.

%%K t16-tit-6 sub
\VUsaut{4.4mm}
\VUpk{10.2pt}{40}{आजोबांची भेट.}
\VUsaut{2.8mm}

%%K t16-17-1 p
१७. --- \VUgras{दातांच्या ब्रशांपैकी} \textsuperscript{(92)} एक मला घेऊन द्यावा अशी मी
आजोबांना विनंती केली — ते ब्रश \VUgras{केसांच्या पिनांशेजारच्या} \textsuperscript{(93)} पेटीत
होते. त्यांनी होकार दिला आणि माझी खरेदी मी स्वतःच गल्ल्यावर जाऊन भरली, जिथे मोठी
\VUgras{रवानगी} \textsuperscript{(100)} तयार केली जात होती, एका \VUgras{गिऱ्हाइकासाठी}
\textsuperscript{(99)}.

%%K t16-18-1 p
१८. --- अचानक कलात्मक \VUgras{ब्राँझच्या कलाकृतींजवळ} \textsuperscript{(101)} थांबलेल्या
माणसांमध्ये थोडी खळबळ उडाली. एक \VUgras{चोर} \textsuperscript{(102)} कोणाचा खिसा कापायचा
प्रयत्न करत असतानाच एका \VUgras{निरीक्षकाने} \textsuperscript{(103)} त्याला रंगेहात पकडले
होते. त्याला \VUgras{अटक} करायला पोलिस बोलवायची व्यवस्था केली गेली, आणि आम्ही बाहेर पडत होतो
तेवढ्यात त्या गुन्हेगाराला तुरुंगात नेले जात होते.

\VUsaut{6.0mm}
\VUfilet{22mm}
